% Bagian: computability
% Bab: recursive-functions
% Subbagian: bounded-minimization

\documentclass[../../../include/open-logic-section]{subfiles}

\begin{document}

\olfileid[id]{cmp}{rec}{bmi}
\olsection{Minimisasi Berbatas}

\begin{explain}
Sering kali berguna untuk mendefinisikan suatu fungsi sebagai bilangan terkecil
yang memenuhi sifat atau relasi~$P$. Jika $P$ dapat diputuskan, kita dapat
menghitung fungsi ini cukup dengan mencoba semua bilangan yang mungkin,
$0$, $1$, $2$, \dots, hingga menemukan bilangan terkecil yang memenuhi~$P$.
Pencarian tak berbatas semacam ini membawa kita keluar dari ranah fungsi
rekursif primitif. Akan tetapi, jika kita hanya tertarik pada bilangan terkecil
yang \emph{kurang dari suatu batas yang diberikan secara terpisah}, kita tetap
berada dalam ranah rekursif primitif. Dengan kata lain, dan secara sedikit lebih
umum, andaikan kita mempunyai relasi rekursif primitif~$R(x,z)$. Tinjau fungsi
yang memetakan $x$ dan~$y$ ke $z < y$ terkecil sedemikian sehingga $R(x, z)$.
Fungsi itu juga dapat dihitung dengan menguji $R(x, 0)$, $R(x, 1)$, \dots,
$R(x, y-1)$. Namun, mengapa fungsi itu bersifat rekursif primitif?
\end{explain}

\begin{prop}
Jika $R(\vec x, z)$ bersifat rekursif primitif, maka demikian pula fungsi
$m_R(\vec{x}, y)$ yang mengembalikan~$z$ terkecil yang kurang dari~$y$
sedemikian sehingga $R(\vec x, z)$ berlaku, jika ada, dan mengembalikan $y$
selain itu. Kita akan menulis fungsi~$m_R$ sebagai
\[
\bmin{z < y}{R(\vec{x}, z)},
\]
\end{prop}

\begin{proof}
Tidak mungkin ada~$z < 0$ sedemikian sehingga $R(\vec{x}, z)$ karena memang
tidak ada $z < 0$. Jadi, $m_R(\vec x, 0) = 0$.

Jika batas berbentuk $y + 1$, terdapat tiga kasus:
\begin{enumerate}
\item Terdapat $z < y$ sedemikian sehingga $R(\vec{x}, z)$, dan dalam kasus
ini $m_R(\vec{x}, y+1) = m_R(\vec{x}, y)$.
\item Tidak terdapat~$z<y$ semacam itu, tetapi $R(\vec{x}, y)$ berlaku; dalam
kasus ini $m_R(\vec{x}, y+1) = y$.
\item Tidak terdapat $z < y+1$ sedemikian sehingga $R(\vec{x}, z)$; dalam
kasus ini $m_R(\vec{x}, y+1) = y+1$.
\end{enumerate}
Jadi, kita dapat mendefinisikan $m_R(\vec x,y)$ dengan rekursi primitif sebagai berikut:
\begin{align*}
m_R(\vec{x}, 0) & = 0\\
m_R(\vec{x}, y+1) & =
\begin{cases}
m_R(\vec{x}, y) & \text{jika $m_R(\vec x, y) \neq y$}\\
y & \text{jika $m_R(\vec x, y) = y$ dan $R(\vec{x}, y)$}\\
y+1 & \text{selain itu.}
\end{cases}
\end{align*}
Perhatikan bahwa terdapat $z<y$ sedemikian sehingga $R(\vec x, z)$ jika dan
hanya jika $m_R(\vec x, y) \neq y$.
\end{proof}

\begin{prob}
Andaikan $R(\vec x, z)$ bersifat rekursif primitif. Definisikan fungsi
$m'_R(\vec{x}, y)$ yang mengembalikan~$z$ terkecil yang kurang dari~$y$
sedemikian sehingga $R(\vec{x}, z)$ berlaku, jika ada, dan mengembalikan $0$
selain itu, dengan rekursi primitif dari~$\Char{R}$.
\end{prob}

\end{document}
